\chapter{Interpretacion de Hallazgos y Patrones de Fatiga}
\label{cap:interpretation}

\section{Signos Fisiologicos de Fatiga Fisica}

Cuando un individuo experimenta fatiga fisica significativa (ejercicio intenso),
se observan patrones predecibles en FatigueSet:

\begin{itemize}
    \item \textbf{HR aumenta dramaticamente}: ~76 bpm (M1) a ~105 bpm (M2), mas incremento con intensidad
    \item \textbf{HRV disminuye}: ~70 ms (M1) a ~49 ms (M3) - el corazon pierde flexibilidad
    \item \textbf{BR se eleva}: respiracion acelerada para suministro de oxigeno
    \item \textbf{EDA sube}: respuesta simpatica a demanda metabolica
    \item \textbf{Autoreporte fisica}: puede aumentar 20-40 puntos en escala
\end{itemize}

La \textbf{disminucion de HRV} es el marcador mas objetivo de fatiga fisica.
En contextos clinicos/laborales, HRV < 30 ms habitualmente indica fatiga severa.

\section{Signos de Fatiga Mental}

Fatiga mental tiene patrones ligeramente diferentes:

\begin{itemize}
    \item \textbf{EEG Theta aumenta}: especialmente en fases posteriores a tareas cognitivas
    \item \textbf{EEG Alpha puede aumentar} si acompanado de somnolencia
    \item \textbf{EDA permanece elevada}: reflejo persistente de demanda atencional
    \item \textbf{HR puede elevarse pero menos que con fatiga fisica}
    \item \textbf{Autoreporte mental}: aumenta posteriormente a tareas (M3)
\end{itemize}

Fatiga mental es mas sutil en senales fisiologicas - EEG es mas informatico que medidas autonomicas.

\section{Recuperacion Observada}

En FatigueSet, la recuperacion es \textbf{incompleta} en los ~10 minutos de M3:

\begin{itemize}
    \item \textbf{HR partially normaliza}: 105 → 84 bpm (~80\% recuperacion)
    \item \textbf{HRV sigue comprometido}: 49 ms aun bajo (recuperacion lenta)
    \item \textbf{EDA persiste elevada}: indica arousale aun presente
    \item \textbf{Autorreportes}: aun moderadamente elevados
\end{itemize}

Esto sugiere que la recuperacion de fatiga es mas lenta que su induccion.
Para fatiga severa, pueden requerirse horas de recuperacion efectiva.

\section{Interacciones Fatiga Fisica × Mental}

En M3 (post-tareas cognitivas despues de actividad fisica), se observan signos
de fatiga combinada:
\begin{itemize}
    \item Ambos autoreportes aumentan
    \item EEG Theta esta elevado (mental)
    \item HRV sigue bajo (residuo de fisica)
\end{itemize}

Aunque los datos no permiten causalidad definitiva, la simultaneidad de fatiga
sugiere competencia de recursos: cuerpo cansado + mente fatigada = estado mas
comprometido que uno solo.

\section{Validacion Cruzada de Modalidades}

Un ejemplo de valor multimodal: Si solo tuvieramos EEG, podriamos confundir
somnolencia con relajacion. Pero combinando:
\begin{itemize}
    \item EEG Alpha/Theta alto
    \item + HR elevado
    \item + EDA elevado
    \item + Autoreporte "cansado mentalmente"
\end{itemize}

Podemos desambiguar: es fatiga (con activacion), no relajacion pasiva.

\section{Implicaciones Clinicas}

\begin{enumerate}
    \item \textbf{Deteccion de fatiga}: Cualquier algoritmo debe usar multiples sensores. HR solo = falsos positivos
    \item \textbf{Individualizacion}: Los umbrales deben normalizarse por participante
    \item \textbf{Tiempo de recuperacion}: Planificar pausas mas largas para recuperacion efectiva
    \item \textbf{Diferenciacion}: Fatiga fisica vs mental requiere EEG para distincion clara
\end{enumerate}

El FatigueSet provide evidencia empirica para todas estas recomendaciones.
